% Long-form recipe -- ctan:xcookybooky, LPPL-1.3 (verified in the shipped
% xcookybooky.sty header: "This file may be distributed and/or modified
% under the conditions of the LaTeX Project Public License, either version
% 1.3 of this license or (at your option) any later version"). Ships with
% TeX Live (tex/latex/xcookybooky) -- starter does not bundle the class. No
% fontspec -- pdflatex-safe.
%
% The package's own \step macro (a \lettrine-numbered preparation step) is
% NOT used here: it passes \stepcounter{step}\thestep -- executable code, not
% a bare letter -- as \lettrine's first argument, and the CURRENT TeX Live
% lettrine (v2.83, dated 2025-12-06; xcookybooky itself dates from 2015) is
% no longer tolerant of that, throwing a hard "Missing \endcsname inserted"
% pile-up on every \end{recipe}. Confirmed by isolating \step down to a bare
% \lettrine[...]{stepcounter...}{}  call with no xcookybooky involved at all
% -- a real upstream incompatibility, not a mistake in this starter's markup.
% Preparation steps below are plain manually-numbered paragraphs instead.
%
% The package's differentiator from the plainer cookybooky/cuisine families
% is genuine support for recipes that run LONG (a wrapped ingredients table
% alongside a preparation column, a \suggestion aside, a \hint callout box),
% so this starter is one deliberately long recipe (a braised short-rib ragu
% with a from-scratch stock sub-recipe) rather than a short card. No
% \graph{} picture is supplied for either recipe -- the key accepts an empty
% value and simply omits the image area, which is the correct choice here
% rather than shipping a placeholder image of unclear provenance.
\documentclass[a4paper]{article}
\usepackage{xcookybooky}

\begin{document}

\begin{recipe}[portion = 6 people, preparationtime = 45 minutes,
  bakingtime = 3 hours, source = House recipe]{Braised Short-Rib Rag\`u}

\graph{}

This is a slow, patient ragu built on braised beef short rib rather than
ground meat -- the long braise breaks the connective tissue down into a
sauce richer than any quick weeknight bolognese. Start the stock the day
before if you can; it keeps well and the ragu is better for it.

\ingredients
{
  2.5\,kg & bone-in beef short ribs \\
  1 & onion, finely diced \\
  2 & carrots, finely diced \\
  2 & celery stalks, finely diced \\
  6 & cloves garlic, minced \\
  200\,ml & red wine \\
  800\,g & crushed tomatoes \\
  1 litre & Brown Beef Stock (see below) \\
  2 & bay leaves \\
  1 & sprig rosemary \\
  3\,tbsp & tomato paste \\
  & olive oil, salt, black pepper \\
  500\,g & pappardelle, to serve \\
}

\preparation
{
\textbf{1.} Season the short ribs generously with salt and pepper. Heat a
heavy casserole over high heat with a splash of olive oil and brown the
ribs on all sides, working in batches so the pan stays hot. Set the ribs
aside.

\textbf{2.} Lower the heat to medium. Add the onion, carrot and celery to
the same pot and cook for 8 minutes until softened. Add the garlic and
tomato paste and cook for 2 minutes more, stirring constantly so the paste
does not catch.

\textbf{3.} Pour in the red wine, scraping up the browned bits from the
bottom of the pot, and simmer for 3 minutes to cook off the alcohol.

\textbf{4.} Return the ribs to the pot. Add the crushed tomatoes, stock, bay
leaves and rosemary. Bring to a bare simmer, cover, and transfer to a
160\textcelsius\ oven for 3 hours, until the meat falls from the bone.

\textbf{5.} Lift the ribs out, discard the bones, and shred the meat back
into the sauce. Skim excess fat from the surface. Simmer uncovered for a
further 15 minutes to thicken. Taste and adjust the seasoning.

\textbf{6.} Cook the pappardelle in well-salted water until al dente, then
toss directly through the ragu with a splash of the pasta water.
}

\suggestion[Make ahead]{The ragu improves overnight in the refrigerator and
freezes well for up to three months -- make a double batch.}

\hint{If the sauce is too thin after 15 minutes of reducing, remove the
lid entirely and simmer harder rather than adding flour or cornstarch: the
gelatine from the bones is doing the thickening work.}

\end{recipe}

\begin{recipe}[portion = about 1.5 litres, preparationtime = 20 minutes,
  bakingtime = 4 hours]{Brown Beef Stock}

\graph{}

A dark, roasted stock is what gives the ragu above its depth. This makes
more than the ragu needs; freeze the rest in small containers.

\ingredients
{
  2\,kg & beef bones (marrow and knuckle) \\
  2 & onions, quartered, skin on \\
  2 & carrots, halved \\
  2 & celery stalks, halved \\
  3\,tbsp & tomato paste \\
  2 & bay leaves \\
  10 & black peppercorns \\
  3 litres & cold water \\
}

\preparation
{
\textbf{1.} Heat the oven to 220\textcelsius. Spread the bones, onions,
carrots and celery on a roasting tray and roast for 45 minutes, turning
once, until deeply browned.

\textbf{2.} Transfer the roasted bones and vegetables to a large stockpot.
Deglaze the roasting tray with a splash of water, scraping up the browned
bits, and add this to the pot along with the tomato paste.

\textbf{3.} Add the cold water, bay leaves and peppercorns. Bring to a
bare simmer -- never a boil -- skimming any foam from the surface for the
first 20 minutes.

\textbf{4.} Simmer uncovered for 4 hours, topping up with water if it
reduces below the level of the bones. Strain through a fine sieve,
discarding the solids, and cool before refrigerating or freezing.
}

\hint{A rolling boil emulsifies the fat into the liquid and clouds the
stock. Keep it at the barest simmer -- the occasional lazy bubble.}

\end{recipe}

\end{document}
